\documentclass[13pt, aspectratio=169]{beamer}

\usepackage{bookmark}
% \usepackage[T1]{fontenc}

% code listings
\usepackage{listings}
\usepackage{xcolor}


% defining colors
\definecolor{uchic_lightgray}{RGB}{214, 214, 206}
\definecolor{uchic_orange}{RGB}{193, 102, 34}
\definecolor{uchic_darkblue}{RGB}{21, 67, 95}

\lstset{
    basicstyle=\ttfamily\small,
    backgroundcolor=\color{gray!10},
    frame=single,
    rulecolor=\color{gray!40},
    breaklines=true,
    showstringspaces=false,
    keywordstyle=\color{uchic_darkblue}\bfseries,
    commentstyle=\color{gray},
    stringstyle=\color{uchic_orange},
}

% beamer theme
\usetheme{Madrid}
\useoutertheme[subsection=false]{miniframes}
\useinnertheme{circles}

% colors
\setbeamercolor{background canvas}{bg=white, fg=lightgray!30}
\setbeamercolor{palette primary}{bg=uchic_darkblue, fg=lightgray!30}
\setbeamercolor{palette secondary}{bg=uchic_lightgray, fg=lightgray!30}
\setbeamercolor{palette tertiary}{bg=uchic_darkblue, fg=lightgray!30}
\setbeamercolor{palette quaternary}{bg=uchic_darkblue, fg=lightgray!30}
\setbeamercolor{structure}{fg=uchic_darkblue}
\setbeamercolor{section in toc}{fg=uchic_darkblue}
\setbeamercolor{title}{bg=uchic_darkblue, fg=lightgray!30}
\setbeamercolor{block title}{bg=uchic_darkblue, fg=white}
\setbeamercolor{block body}{bg=gray!10}
\setbeamercolor{alerted text}{fg=uchic_orange}
\setbeamercolor{section in head/foot}{bg=uchic_darkblue, fg=white}
\setbeamercolor{mini frame}{bg=uchic_darkblue, fg=white}

\usepackage{etoolbox}
\patchcmd{\sectionentry}{\usebeamercolor[fg]{section in head/foot}}{\usebeamercolor[fg]{mini frame}}{}{}

\setbeamercovered{transparent}
\setbeamertemplate{footline}[frame number]
\setcounter{subsection}{1}
\setbeamertemplate{navigation symbols}{}

% title page
\setbeamertemplate{title page}{
  \vbox{}
  \vfill
  \begingroup
    \centering
    \usebeamertemplate{title}
    \vskip0.5em
    \usebeamertemplate{titlegraphic}
    \vskip1em
    \usebeamertemplate{author}
    \begin{center}
        \vfill
        \footnotesize{Economic Society for Bocconi Students}
    \end{center}
    \footnotesize{\usebeamertemplate{date}}
  \endgroup
}

\title{Good Coding Practices in Economics}
\author{Jakub Przewoski}
\date{March 2026}

% ============================================================
\begin{document}

\maketitle

% ---- TABLE OF CONTENTS ----
\begin{frame}{Today's Agenda}
    \tableofcontents
\end{frame}

% ============================================================
\section{Introduction}

\begin{frame}{Why Does This Matter?}
    \begin{columns}
        \begin{column}{0.55\textwidth}
            \begin{itemize}
                \item Code is \textbf{mostly read, not written} — by colleagues, supervisors, and your future self
                \item Open-source languages and conventions are a long-term investment in efficiency
                \item Today: a quick overview of core concepts + a \alert{live demo}
            \end{itemize}
        \end{column}
        \begin{column}{0.45\textwidth}
            \begin{block}{The Golden Rule}
                Always code as if the person who will maintain your code is a violent psychopath who knows where you live.
            \end{block}
        \end{column}
    \end{columns}
\end{frame}

% ============================================================
\begin{frame}{What Is Reproducibility?}
    \textbf{Scenario:} A professor asks you to write code he can run on his machine in one click. Later, he'll adjust it and eventually take over the project entirely.

    \vspace{0.5em}
    What should your code guarantee?

    \begin{itemize}
        \item \textbf{Dependency transparency} — what libraries does it need, and which versions?
        \item \textbf{Documentation} — can someone understand it quickly without asking you?
        \item \textbf{Modularity} — can individual parts be changed without breaking everything else?
    \end{itemize}

    \vspace{0.4em}
    \begin{alertblock}{The core goal}
        Anyone with your code and your instructions should get \textit{exactly} the same results.
    \end{alertblock}
\end{frame}

% ============================================================
\section{Setup}

\begin{frame}{R, Python, or Stata?}
    \begin{columns}
        \begin{column}{0.3\textwidth}
            \textbf{Stata}
            \begin{itemize}
                \item Easy to learn
                \item Limited scope
                \item Requires a license
                \item Rare outside academia
            \end{itemize}
        \end{column}
        \begin{column}{0.3\textwidth}
            \textbf{R}
            \begin{itemize}
                \item Excellent for statistics
                \item Rich econometrics ecosystem
                \item Less general-purpose
                \item Free \& open source
            \end{itemize}
        \end{column}
        \begin{column}{0.3\textwidth}
            \textbf{Python}
            \begin{itemize}
                \item General purpose
                \item ML, simulation, data
                \item Stat inference improving
                \item Free \& widely used
            \end{itemize}
        \end{column}
    \end{columns}

    \vspace{0.8em}
    \begin{block}{Disclosure}
        Today we will be working with Python, but the lessons here are language-agnostic. IRL, it's probably best if you learn a bit of all three.
    \end{block}
\end{frame}

% ============================================================
\begin{frame}{The Tools You Need}
    \begin{columns}
        \begin{column}{0.48\textwidth}
            \textbf{Language \& Package Manager}
            \begin{itemize}
                \item \textbf{Miniconda} — lightweight, terminal-based (recommended)
                \item \textbf{Anaconda} — includes a GUI, easier to start with
                \item \textbf{R} from CRAN:\@ \texttt{cran.r-project.org}
            \end{itemize}
        \end{column}
        \begin{column}{0.48\textwidth}
            \textbf{Code Editor}
            \begin{itemize}
                \item \textbf{VS Code} — lightweight, highly customizable, free
                \item Open it through the terminal or project folder to set up a workspace
            \end{itemize}
        \end{column}
    \end{columns}
\end{frame}

\begin{frame}[fragile]{Conda Environments}
    \textbf{The problem:} Different projects may need different (sometimes conflicting) package versions.

    \vspace{0.3em}
    \textbf{The solution:} A \alert{conda environment} is an isolated, self-contained Python installation per project.

    \begin{columns}
        \begin{column}{0.48\textwidth}
            \begin{itemize}
                \item Each project gets its own packages and versions
                \item Environments don't interfere with each other
                \item Easy to share: export to \texttt{environment.yml}
                \item \textbf{R alternative}: The \texttt{renv} package
            \end{itemize}
        \end{column}
        \begin{column}{0.48\textwidth}
\begin{lstlisting}[language=bash]
# create
conda create -n myproject python=3.11
# activate
conda activate myproject
# export
conda env export > environment.yml
\end{lstlisting}
        \end{column}
    \end{columns}
\end{frame}

\begin{frame}{VS Code Extensions}
    % using VS Code is nicer because more versatile than IDEs and more lightweight but use whatever you prefer!
    Activate extensions \textit{per workspace} to keep startup time fast and battery usage low.

    \vspace{0.5em}
    \begin{columns}
        \begin{column}{0.48\textwidth}
            \textbf{Must-haves}
            \begin{itemize}
                \item \textbf{Python} — linting, autocomplete, IntelliSense
                \item \textbf{Jupyter} — run notebooks and interactive \texttt{.py} mode inside VS Code
                \item \textbf{Data Wrangler} — visual data viewer
            \end{itemize}
        \end{column}
        \begin{column}{0.48\textwidth}
            \textbf{Nice to have}
            \begin{itemize}
                \item \textbf{GitHub Copilot} — free with \href{https://github.com/education}{GitHub Education}
                \item \textbf{LaTeX Workshop} — write papers locally inside VS Code
                \item \textbf{R extension} — if you use R
            \end{itemize}
        \end{column}
    \end{columns}
\end{frame}

% ============================================================
\section{Writing Good Code}

\begin{frame}[fragile]{A Clean Project Folder}
    Good structure means anyone (including future-you) can navigate the project in seconds.

    \vspace{0.5em}
    \begin{columns}
        \begin{column}{0.45\textwidth}
\begin{lstlisting}
my_project/
|-- data/
|   |-- raw/
|   |-- clean/
|-- src/
|   |-- 01_clean.py
|   |-- 02_analysis.py
|   |-- __main__.py
|-- output/
|   |-- figures/
|   |-- tables/
|-- writing/
|-- README.md
|-- environment.yml
\end{lstlisting}
        \end{column}
        \begin{column}{0.5\textwidth}
            \begin{itemize}
                \item \textbf{Raw data is never modified} — always save a clean copy separately
                \item \textbf{README.md} is not optional — it's the front door to your project
                \item Keep one \textbf{sample project folder} and copy it every time you start something new
            \end{itemize}
        \end{column}
    \end{columns}
\end{frame}

% ============================================================
\begin{frame}{The Essentials}
    \begin{itemize}
        \item \textbf{Separate files for separate tasks} — cleaning, analysis, and a \texttt{\_\_main\_\_.py} that ties it together
        \item \textbf{DRY — Don't Repeat Yourself} — if you copy-paste code more than twice, write a function
        \item \textbf{Comment the \textit{why}, not the \textit{what}} — the code already says what it does; explain your reasoning
        \item \textbf{Write docstrings} — a few lines above each function describing inputs, outputs, and purpose
        \item \textbf{Name variables consistently} — \texttt{gdp\_growth\_2023}, not \texttt{x2} or \texttt{temp\_var\_final\_v3}
        \item \textbf{Save intermediate clean data} — avoid re-running expensive cleaning steps on every run
    \end{itemize}
\end{frame}

\begin{frame}[fragile]{Comments \& Docstrings}
    \begin{columns}
        \begin{column}{0.48\textwidth}
            \textbf{Bad comment — states the obvious:}
\begin{lstlisting}[language=python]
# multiply x by 2
result = x * 2
\end{lstlisting}
            \vspace{0.5em}
            \textbf{Good comment — explains the why:}
\begin{lstlisting}[language=python]
# Annualize monthly growth rate
# using compound formula
result = (1 + x)**12 - 1
\end{lstlisting}
        \end{column}
        \begin{column}{0.48\textwidth}
            \textbf{A docstring:}
\begin{lstlisting}[language=python]
def clean_cpi(df, base_year):
    """
    Rebase CPI series to a given
    base year.

    Parameters:
        df: pd.DataFrame with 'date'
            and 'cpi' columns
        base_year: int

    Returns:
        pd.DataFrame with rebased CPI
    """
\end{lstlisting}
        \end{column}
    \end{columns}
\end{frame}

\begin{frame}{Useful Tricks}
    \begin{columns}
        \begin{column}{0.48\textwidth}
            \textbf{Interactive mode in VS Code}
            \begin{itemize}
                \item Use \texttt{Shift + Enter} to run your \texttt{.py} file like a notebook (run code block-by-block) 
                \item Great for exploratory analysis while keeping your code in proper scripts
            \end{itemize}
        \end{column}
        \begin{column}{0.48\textwidth}
            \textbf{Logging \& resumable pipelines}
            \begin{itemize}
                % chktex-file 36 
                \item When developing the code, use the \texttt{print()} statements to test it, 
                \item Before running the code, switch to Python's \texttt{logging} module writes different statements (e.g.\ errors) to a file
                % \item For long-running scripts (e.g.\ web scraping, simulations), design them so they can be stopped and restarted without re-doing completed work
            \end{itemize}
        \end{column}
    \end{columns}

    \vspace{0.5em}
    \begin{block}{AI tools}
        GitHub Copilot is excellent for boilerplate and docstrings. Use it as a \textit{first draft} — always read and understand what it writes (or use Claude Code instead).
    \end{block}
\end{frame}

% ============================================================
\section{Git}

\begin{frame}{What Is Git?}
    \textbf{Git} is a version control system — it tracks every change you make to your code over time.

    \vspace{0.5em}
    \begin{columns}
        \begin{column}{0.55\textwidth}
            \textbf{What it lets you do:}
            \begin{itemize}
                \item See exactly what changed between versions and when
                \item Revert to any previous state if something breaks
                \item Experiment safely using \textbf{branches} — a separate line of development that won't touch your working code
                \item Collaborate with others and \textbf{merge} changes together
                \item See who wrote which line and why (commit messages)
            \end{itemize}
        \end{column}
        \begin{column}{0.4\textwidth}
            \begin{alertblock}{Think of it as \dots}
                Track Changes in Word, but for your entire project — and you control exactly when a snapshot is taken.
            \end{alertblock}
        \end{column}
    \end{columns}
\end{frame}

\begin{frame}[fragile]{Git: The Basic Workflow}
    \begin{columns}
        \begin{column}{0.48\textwidth}
\begin{lstlisting}[language=bash]
# Start tracking a project
git init

# Stage your changes
git add .

# Save a snapshot
git commit -m "Add CPI cleaning script"

# See what changed
git diff

# See your history
git log --oneline
\end{lstlisting}
        \end{column}
        \begin{column}{0.48\textwidth}
            \begin{itemize}
                \item \textbf{GitHub} hosts your repository online — backup + collaboration
                \item \textbf{GitHub Education} gives you free Pro access (private repos, Copilot, and more) with your university email
                \item Always include a \texttt{.gitignore} to avoid committing data files, credentials, or large outputs
            \end{itemize}
        \end{column}
    \end{columns}
\end{frame}

\begin{frame}[fragile]{Branches: Experiment Safely}
    \begin{block}{The idea}
        A branch is a parallel version of your project. Make changes there freely — your main code stays untouched until you decide to merge.
    \end{block}

    \vspace{0.5em}
    \begin{columns}
        \begin{column}{0.48\textwidth}
\begin{lstlisting}[language=bash]
# Create and switch to a new branch
git checkout -b new-specification

# ... make changes, commit ...

# Merge back when ready
git checkout main
git merge new-specification
\end{lstlisting}
        \end{column}
        \begin{column}{0.48\textwidth}
            \textbf{Useful for:}
            \begin{itemize}
                \item Testing a different model specification
                \item Trying a different data cleaning approach
                \item Responding to reviewer comments without touching working code
            \end{itemize}
        \end{column}
    \end{columns}
\end{frame}

\section*{Takeaways}

\begin{frame}{Takeaways}
    \begin{enumerate}
        \item \textbf{Reproducibility} is not optional in research — document dependencies, structure, and reasoning
        \item \textbf{Python and R} are worth learning; they'll serve you in academia and industry
        \item \textbf{Environments} keep projects isolated and shareable
        \item \textbf{Structure} your project before writing a single line of code
        \item \textbf{Write for readers} — comment the why, use functions, keep it DRY
        \item \textbf{Use Git} — you will eventually break something, and you'll be glad you did
    \end{enumerate}
\end{frame}

\begin{frame}{Further Reading}
    \begin{itemize}
        \item \textbf{Missing Semester} (MIT) — \texttt{missing.csail.mit.edu} \\
        Practical tools every developer should know: shell, git, editors, and more
        \vspace{0.3em}
        \item \textbf{Coding for Economists} — \texttt{aeturrell.github.io/coding-for-economists} \\
        Reproducible analysis and best practices tailored for economics research
        \vspace{0.3em}
        \item \textbf{Why Every Economics Student Should Use GitHub} \\
        \texttt{econ-jobs.com} — motivation and practical starting points
    \end{itemize}
    
\end{frame}

% ---

\appendix

\section{Appendix: Local \LaTeX}

\begin{frame}{Writing Papers Locally with \LaTeX}
    \textbf{Why bother with local \LaTeX\ instead of Overleaf?}

    \begin{columns}
        \begin{column}{0.48\textwidth}
            \textbf{Advantages}
            \begin{itemize}
                \item Figures and output live in the same project --- no manual transfers
                \item No compile limits, full privacy, works offline
                \item Integrates with GitHub for version control of your writing
                \item Zotero integration auto-updates your bibliography
                \item AI completions via Copilot
            \end{itemize}
        \end{column}
        \begin{column}{0.48\textwidth}
            \textbf{Disadvantages}
            \begin{itemize}
                \item Initial setup takes some effort
                \item Debugging (e.g.\ bibliography issues) requires more patience
                \item Real-time collaboration is less smooth than Overleaf
            \end{itemize}
            \vspace{0.5em}
            \textbf{VS Code extensions:}
            \begin{itemize}
                \item \textbf{LaTeX} + \textbf{LaTeX Workshop}
            \end{itemize}
        \end{column}
    \end{columns}
\end{frame}

% ============================================================
\end{document}
